\chapter{Introducción y contexto}
\label{ch:introduccion}

% Real draft grounded on the owner's approved TFG proposal. No invented
% citations: bibliographic support that still needs a source is marked with
% \pendiente{citar: ...}.

\section{Planteamiento del problema}
\label{sec:planteamiento}

En la actualidad, el uso de datos en el ámbito deportivo ha adquirido una
relevancia creciente para mejorar el rendimiento, prevenir lesiones y optimizar
los entrenamientos. El presente Trabajo Fin de Grado propone el desarrollo de un
dashboard inteligente orientado al análisis del rendimiento de corredores,
integrando métricas derivadas de fuentes de datos deportivas como el ritmo, la
frecuencia cardíaca, la distancia, la potencia o la cadencia. El proyecto
pretende ofrecer una herramienta de apoyo a la toma de decisiones basada en
datos, capaz de identificar patrones de evolución y de proponer mejoras en los
planes de entrenamiento.

El problema que motiva el trabajo puede formularse así: el corredor y su
entrenador disponen hoy de una gran cantidad de datos, pero no siempre de los
medios para interpretarlos de forma personalizada y sistemática. Convertir ese
volumen de registros en decisiones de entrenamiento fundamentadas es,
precisamente, el objetivo del proyecto.

\section{Contexto}
\label{sec:contexto}

El auge de los dispositivos \emph{wearables} (Garmin, Polar, Suunto) y de
plataformas como Strava o TrainingPeaks ha impulsado la generación masiva de
datos de rendimiento deportivo. Sin embargo, la mayoría de estos sistemas
presentan interfaces cerradas o de uso limitado para el análisis personalizado.
El desarrollo de dashboards a medida mediante herramientas como Power BI, Python
(Dash, Streamlit) o React permite crear soluciones adaptadas a las necesidades
del usuario final, integrando múltiples fuentes y técnicas de análisis
predictivo.

Este trabajo se enmarca dentro de las metodologías de \emph{Data-Driven
Management} aplicadas al contexto deportivo, esto es, el uso sistemático de
datos y de análisis cuantitativo para respaldar decisiones
\cite{provost2013,davenport2010}. La generalización del registro automático de
actividad ha hecho que ese enfoque sea hoy viable también para deportistas
individuales y no solo para organizaciones \cite{marr2016}.

\pendiente{citar: estudios revisados por pares sobre el uso de wearables y
plataformas deportivas en poblaciones de corredores aficionados y de
competición.}

\section{Motivación}
\label{sec:motivacion}

La motivación del trabajo es doble. Por un lado, la cantidad y la variedad de
métricas disponibles (ritmo, frecuencia cardíaca, potencia, cadencia, distancia)
permiten caracterizar la evolución del corredor con un detalle que antes no era
accesible. Por otro lado, el tratamiento estadístico y el aprendizaje automático
ofrecen herramientas consolidadas para extraer patrones y para construir
predicciones a partir de esos registros \cite{montgomery2018,geron2023}.

La oportunidad, por tanto, no está en generar más datos, sino en integrarlos y
presentarlos de forma que resulten accionables tanto para un corredor como para
un entrenador que gestiona varios atletas.

\pendiente{citar: literatura específica sobre análisis de rendimiento en carrera
y sobre el paso de la métrica aislada a la decisión de entrenamiento.}

\section{Relevancia}
\label{sec:relevancia}

La relevancia del proyecto se apoya en tres ideas. En primer lugar, existe una
laguna entre los datos que ya se registran y su aprovechamiento efectivo para la
toma de decisiones. En segundo lugar, las soluciones actuales tienden a estar
cerradas o limitadas al análisis personalizado, de modo que una herramienta
adaptable cubre una necesidad real. Por último, la incorporación de un módulo
analítico y predictivo permite pasar de la descripción retrospectiva a la
estimación de la progresión del rendimiento, lo que aporta valor tanto al
deportista individual como al entrenador o al equipo.

\pendiente{citar: trabajos que documenten la brecha entre recogida de datos y
explotación en plataformas deportivas comerciales, para sostener la afirmación
anterior con evidencia.}
